% !TeX root = ../../Book2.tex
% 纲要标题：## 第4章　系数变换、乘积与上同调环
% 纲要位置：几何学计划纲要.md，第 1151 行。
% 正文状态：初稿；尚未正式审核。

% 纲要内容开始
% 
% > 本章目标与证明范围：在确实需要系数变化时建立自由消解、Ext 与 Tor。随后从链级乘积推出环结构和 Künneth；不把自然正合列的非典范分裂写成自然同构。
% 
% 纲要内容结束

\chapter{系数变换、乘积与上同调环}
\label{g2:ch:04}
\BookChapterNumber{g2:ch:04}

\begin{chapterabstract}
	本章研究系数与乘法怎样丰富空间的不变量。自由消解把取同态、取张量时的正合性障碍组织为 Ext 与 Tor，万有系数定理据此从自由链复形的同调恢复任意系数同调及上同调。Alexander--Whitney 映射与 shuffle 映射比较乘积空间的链与张量复形，导出 Künneth 短正合列。随后通过对角映射定义杯积，证明自然性、结合性、单位与分次交换性，并建立相对乘积和取值配对。环面、闭曲面及复射影空间的计算说明，即使加性同调相同，乘法仍能区分不同的同伦类型。
\end{chapterabstract}

\begin{learninggoals}
	\begin{itemize}
		\item 构造自由消解及其比较链映射，计算循环模的 Ext、Tor，并说明两者分别测量哪种正合性障碍。
		\item 使用两种万有系数正合列计算系数同调与上同调，区别自然映射和依赖选择的分裂。
		\item 写出 Alexander--Whitney 与 shuffle 映射，核对总微分及交换因子的符号，用 Künneth 定理计算积空间中的自由类与挠类。
		\item 从上链定义杯积，证明其环性质及分次交换性，判断相对杯积的目标及所需切除条件。
		\item 计算环面、可定向曲面、不可定向曲面的模二环及复射影空间的上同调环，用乘法关系约束映射和区分同伦类型。
	\end{itemize}
\end{learninggoals}

上一章中，射影平面的胞腔链含有一条乘以 \(2\) 的边界。用整数计算，它产生一次同调中的二阶挠类，却不产生二次同调；把系数改成 \(\mathbb F_2\) 后，这条边界成为零，二次圈反而出现了。系数变化不仅改变已有群中元素的写法，还可能把原来有边界的链变成新的圈。因此，需要在链复形上处理系数变化，再计算其同调。

\ProseParagraphBreak

另一个问题来自 \(\mathbb CP^2\) 与 \(S^2\vee S^4\)。两者的整数同调都在次数 \(0,2,4\) 各有一个 \(\mathbb Z\)，单看这些群无法区分它们。可是 \(\mathbb CP^2\) 中两条坐标射影直线交于一点，这暗示两个二次信息同时出现时可能留下一个四次信息；在球面的楔和中，两个维数的部分则只在基点接触。杯积会把这种差别变成可以证明、可以计算的乘法关系。

\ProseParagraphBreak

系数变化带来的新圈和新上同调类，将由自由消解、Ext 与 Tor 描述；这些构造从模、核、商、自由基和线性映射开始建立。整数与域上的计算是主要例子，采用主理想整环的统一表述，是因为其自由模的子模仍然自由，足以把所需障碍放进两项消解。转向乘积后，链的张量构造把两组维数合为总次数；Künneth 定理据此计算乘积同调，杯积再沿对角映射给出同一空间上的乘法。杯积只要求系数为含幺交换环，Künneth 的短正合列则使用主理想整环条件。

\ProseParagraphBreak

几何部分使用\chapref{g2:ch:02}的奇异链、链同伦、空间对正合列和切除：链同伦用于比较乘积的不同链模型，正合列与切除用于构造相对乘积。\chapref{g2:ch:03}的胞腔计算则提供换系数前的整数同调。关于上同调与万有系数、乘积与杯积的阅读对照，可分别查\cite[第3.1节、第3.A节；第3.2节、第3.B节]{Hatcher2002}。

% !TeX root = ../../Book2.tex
% 纲要标题：### 4.1 自由消解与系数障碍
% 纲要位置：几何学计划纲要.md，第 1105 行。
% 正文状态：未撰写；本文件目前是框架。

% 纲要内容开始
% 
% * 自由消解和比较同伦
% * Ext、Tor 的定义与整数循环模计算
% * 维数平移的最低形式
% * 张量复形与符号约定
% 
% 纲要内容结束

\section{自由消解与系数障碍}
\label{g2:sec:0401}

待撰写

\subsection{自由消解和比较同伦}
\label{g2:subsec:0401:01}
待撰写

\subsection{Ext、Tor 的定义与整数循环模计算}
\label{g2:subsec:0401:02}
待撰写

\subsection{维数平移的最低形式}
\label{g2:subsec:0401:03}
待撰写

\subsection{张量复形与符号约定}
\label{g2:subsec:0401:04}
待撰写
% 纲要条目排为编号小节标题，并进入目录与书签；说明段仍保留于纲要注释中。

% !TeX root = ../../Book2.tex
% 纲要标题：### 4.2 万有系数定理
% 纲要位置：几何学计划纲要.md，第 1112 行。
% 正文状态：未撰写；本文件目前是框架。

% 纲要内容开始
% 
% * 自由链复形的圈边界分解
% * 上同调的自然短正合列
% * 同调的系数变换形式
% * 非典范分裂及函子性的限制
% 
% 纲要内容结束

\section{万有系数定理}
\label{g2:sec:0402}

待撰写

\subsection{自由链复形的圈边界分解}
\label{g2:subsec:0402:01}
待撰写

\subsection{上同调的自然短正合列}
\label{g2:subsec:0402:02}
待撰写

\subsection{同调的系数变换形式}
\label{g2:subsec:0402:03}
待撰写

\subsection{非典范分裂及函子性的限制}
\label{g2:subsec:0402:04}
待撰写
% 纲要条目排为编号小节标题，并进入目录与书签；说明段仍保留于纲要注释中。

% !TeX root = ../../Book2.tex
% 纲要标题：### 4.3 空间乘积与 Eilenberg–Zilber
% 纲要位置：几何学计划纲要.md，第 1119 行。
% 正文状态：未撰写；本文件目前是框架。

% 纲要内容开始
% 
% * Alexander–Whitney 与 shuffle 映射
% * 链同伦等价和符号
% * 域及主理想整环上的 Künneth 定理
% * Tor 项和积空间算例
% 
% 纲要内容结束

\section{空间乘积与 Eilenberg–Zilber}
\label{g2:sec:0403}

待撰写

\subsection{Alexander–Whitney 与 shuffle 映射}
\label{g2:subsec:0403:01}
待撰写

\subsection{链同伦等价和符号}
\label{g2:subsec:0403:02}
待撰写

\subsection{域及主理想整环上的 Künneth 定理}
\label{g2:subsec:0403:03}
待撰写

\subsection{Tor 项和积空间算例}
\label{g2:subsec:0403:04}
待撰写
% 纲要条目排为编号小节标题，并进入目录与书签；说明段仍保留于纲要注释中。

% !TeX root = ../../Book2.tex
% 纲要标题：### 4.4 杯积与环结构
% 纲要位置：几何学计划纲要.md，第 1176 行。
% 正文状态：初稿；尚未正式审核。

% 纲要内容开始
% 
% * 对角映射及杯积
% * 自然性、结合性与单位
% * 分次交换性的链同伦证明
% * 相对杯积及配对
% 
% 纲要内容结束

\section{杯积与环结构}
\label{g2:sec:0404}

同调外积把两个空间的类放到乘积空间中。如果两个类原来就来自同一个空间，还需要一条把它们放回该空间的自然映射。对角映射 \(x\mapsto(x,x)\) 正好提供这条映射，但它的方向是从 \(X\) 到 \(X\times X\)。因此，适合沿它拉回的对象是上同调。杯积把“同一点的两份信息同时取值”转成乘法，使上同调能够记录加性群看不见的关系。本节使用任意含幺交换环 \(R\)，上边界始终取 \(\delta\varphi=\varphi\partial\)。有关杯积及其相对形式可参阅 \cite[第3.2节]{Hatcher2002}。

\subsection{对角映射及杯积}
\label{g2:subsec:0404:01}

先把两个上链作用在张量链上的含义说清。若 \(\varphi\in C^p(X;R)\)、\(\psi\in C^q(Y;R)\)，定义次数 \(p+q\) 的线性泛函
\[
	(\varphi\otimes\psi)(a\otimes b)=\varphi(a)\psi(b)
	\quad(|a|=p,\ |b|=q),
\]
并在其他双次数上取零。这只是一个明确的泛函定义，不要求
\(\operatorname{Hom}(C_p,R)\otimes\operatorname{Hom}(C_q,R)\to
\operatorname{Hom}(C_p\otimes C_q,R)\) 为同构。按上一节的总微分，
\begin{equation}\label{g2:eq:tensor-cochain-coboundary}
	\delta(\varphi\otimes\psi)
	=\delta\varphi\otimes\psi+(-1)^p\varphi\otimes\delta\psi.
\end{equation}
所需符号由张量总微分给出。

\begin{definition}[杯积]\label{g2:def:cup-product}
	设 \(X\) 为拓扑空间，\(R\) 为含幺交换环。令 \(\Delta_X:X\to X\times X\) 为对角映射，记
	\[
		D_X=\operatorname{AW}\circ(\Delta_X)_\#
		:C_*(X;R)\longrightarrow C_*(X;R)\otimes_R C_*(X;R).
	\]
	对 \(\varphi\in C^p(X;R)\)、\(\psi\in C^q(X;R)\)，定义其\term[beiji]{杯积}[cup product]为
	\[
		\varphi\smile\psi=(\varphi\otimes\psi)D_X.
	\]
	等价地，对 \(\sigma:\Delta^{p+q}\to X\)，
	\begin{equation}\label{g2:eq:cup-front-back}
		(\varphi\smile\psi)(\sigma)
		=\varphi\bigl(\sigma[0,\ldots,p]\bigr)
			\psi\bigl(\sigma[p,\ldots,p+q]\bigr).
	\end{equation}
\end{definition}

切分点的位置由两个上链的次数决定：前 \(p\) 维面给第一个上链，后 \(q\) 维面给第二个。三角形上的一次上链杯积只使用边 \(01\) 与边 \(12\)；边 \(02\) 不直接出现在这条取值公式中，却通过上链的闭性参与不同代表之间的比较。这种不对称属于链层面的选择，稍后会证明它在上同调中只留下分次交换符号。\symidx{\(\smile\)：上链与上同调的杯积}

\begin{proposition}[杯积的上边界]\label{g2:prop:cup-leibniz}
	对拓扑空间 \(X\)、含幺交换环 \(R\) 以及 \(\varphi\in C^p(X;R)\)、\(\psi\in C^q(X;R)\)，有
	\[
		\delta(\varphi\smile\psi)
		=\delta\varphi\smile\psi+(-1)^p\varphi\smile\delta\psi.
	\]
	因此杯积诱导双线性映射
	\[
		H^p(X;R)\otimes_R H^q(X;R)\longrightarrow H^{p+q}(X;R),
		\qquad
		[\varphi]\otimes[\psi]\longmapsto[\varphi\smile\psi].
	\]
\end{proposition}
\begin{proof}
	由\cref{g2:prop:aw-shuffle-chain-maps}，\(D_X\) 是链映射。将\eqref{g2:eq:tensor-cochain-coboundary}与 \(D_X\) 复合，即得所述等式。如果 \(\delta\varphi=\delta\psi=0\)，其杯积也闭；若将 \(\varphi\) 改成 \(\varphi+\delta u\)，杯积的变化为 \(\delta(u\smile\psi)\)。若将 \(\psi\) 改成 \(\psi+\delta v\)，变化为 \((-1)^p\delta(\varphi\smile v)\)。所以所得上同调类只依赖两个原类。
\end{proof}

\subsection{自然性、结合性与单位}
\label{g2:subsec:0404:02}

杯积的边界公式使它下降到上同调。接下来比较拉回、三次相乘和常值零次上链与这项运算的关系：拉回保持单形的各个面，三次相乘按两个切分点分割单形，常值 \(1\) 则在零维端面上给出单位元。

\begin{theorem}[杯积的环性质]\label{g2:thm:cup-properties}
	设 \(R\) 为含幺交换环，\(X\) 为拓扑空间。杯积使
	\[
		H^*(X;R)\coloneqq\bigoplus_{n\ge0}H^n(X;R)
	\]
	成为含幺分次 \(R\) 代数，乘法把次数 \(p,q\) 相加。其单位是把每个点取值为 \(1\) 的零次上圈所代表的类。对连续映射 \(f:X\to Y\)，
	\[
		f^*(\alpha\smile\beta)=f^*\alpha\smile f^*\beta,
		\qquad f^*(1_Y)=1_X.
	\]
	特别地，同伦等价诱导上同调环同构。
\end{theorem}
\begin{proof}
	双线性已经证明。若 \(\varphi,\psi,\eta\) 的次数分别为 \(p,q,r\)，则两种结合方式在任意 \((p+q+r)\) 单形 \(\sigma\) 上都给出
	\[
		\varphi(\sigma[0,\ldots,p])\,
		\psi(\sigma[p,\ldots,p+q])\,
		\eta(\sigma[p+q,\ldots,p+q+r]).
	\]
	因此结合律在上链上就严格成立。令 \(\mathbf1\) 在每个零单形上取值 \(1\)，则 \((\delta\mathbf1)(\gamma)=1-1=0\)；前后面公式分别取切分位置 \(0\) 和 \(p\)，直接得到
	\[
		\mathbf1\smile\varphi=\varphi=\varphi\smile\mathbf1.
	\]
	这给出单位。空空间的上同调为零，按其零环的通常约定处理。

	连续映射与取面交换，故
	\[
		\bigl(f^*(\varphi\smile\psi)\bigr)(\sigma)
		=\varphi(f\sigma[0,\ldots,p])\psi(f\sigma[p,\ldots,p+q])
		=\bigl(f^*\varphi\smile f^*\psi\bigr)(\sigma).
	\]
	单位也保持。最后，\cref{g2:thm:singular-homotopy-invariance}的棱柱链同伦经 \(\operatorname{Hom}(-,R)\) 给出上链同伦：若 \(g_\#-f_\#=\partial P+P\partial\)，则对上圈 \(\varphi\) 有 \(g^*\varphi-f^*\varphi=\delta(\varphi P)\)。故同伦映射诱导相同的上同调映射；把它用于同伦等价及其逆，即得环同构。
\end{proof}

我们称这个分次代数为 \(X\) 的\term[shangtongdiao huan]{上同调环}[cohomology ring]。零次上同调可以理解为每个道路分支上的常值函数，杯积在零次就是逐点相乘。若 \(X\) 道路连通，则 \(H^0(X;R)=R\)，其中的元素在其他次数上按通常的标量乘法作用。

\subsection{分次交换性的链同伦证明}
\label{g2:subsec:0404:03}

一般上链的杯积并不交换：在一个三角形上，先看边 \(01\) 再看边 \(12\)，与把两条边的任务互换，通常给出不同的数。要证明上同调上的交换律，必须找到两种取值之差的原上链，而不能直接交换前后面。

\begin{proposition}[分次交换性]\label{g2:prop:graded-commutativity}
	设 \(X\) 为拓扑空间，\(R\) 为含幺交换环。对 \(\alpha\in H^p(X;R)\)、\(\beta\in H^q(X;R)\)，
	\[
		\alpha\smile\beta=(-1)^{pq}\beta\smile\alpha.
	\]
	更具体地，存在自然算子 \(K:C_n(X;R)\to(C_*(X;R)\otimes C_*(X;R))_{n+1}\)，使
	\begin{equation}\label{g2:eq:diagonal-switch-homotopy}
		\tau D_X-D_X=\mathrm{d}K+K\partial,
		\qquad
		\tau(a\otimes b)=(-1)^{|a||b|}b\otimes a.
	\end{equation}
\end{proposition}
\begin{proof}
	先核对 \(\tau\) 是链映射。对 \(|a|=i\)、\(|b|=j\)，\(\mathrm{d}\tau(a\otimes b)\) 的两项系数分别为 \((-1)^{ij}\)、\((-1)^{ij+j}\)；\(\tau\mathrm{d}(a\otimes b)\) 中对应的 \(b\otimes\partial a\)、\(\partial b\otimes a\) 项系数分别为 \((-1)^{(i-1)j}\)、\((-1)^{i+i(j-1)}\)。前者与 \(ij+j\)、后者与 \(ij\) 模二相同，故两式一致。

	于是 \(D_X\) 与 \(\tau D_X\) 是两条自然链映射。它们在零次均把一个点送到该点与自身的张量。源的通用基元是 \(\operatorname{id}_{\Delta^n}\)，目标模型 \(C_*(\Delta^n;R)\otimes C_*(\Delta^n;R)\) 的增广收缩已在\cref{g2:lem:product-model-homotopy}后给出。应用该引理便得到 \(K\)。具体地，已知低维 \(K\) 后，在模型中填充圈
	\[
		(\tau D-D)\operatorname{id}_{\Delta^n}
			-K\partial\operatorname{id}_{\Delta^n},
	\]
	再沿 \(\sigma\otimes\sigma\) 推到任意奇异单形 \(\sigma\)。这给出所需的自然递归构造，而且 \(K\sigma\) 的两个因子始终留在 \(\sigma(\Delta^n)\) 内。

	取代表上圈 \(\varphi,\psi\)，并令 \(\lambda=\varphi\otimes\psi\)。由\eqref{g2:eq:tensor-cochain-coboundary}，\(\lambda\mathrm{d}=0\)。把 \(\lambda\) 作用在\eqref{g2:eq:diagonal-switch-homotopy}两边，得到
	\[
		(-1)^{pq}\psi\smile\varphi-\varphi\smile\psi
		=(\lambda K)\partial=\delta(\lambda K).
	\]
	故差为上边界，所述上同调等式成立。若 \(p=q=0\)，右侧没有负次上链，此时两种取值已由 \(R\) 的交换律严格相同。
\end{proof}

这样的环称为\term[fenci jiaohuan huan]{分次交换环}[graded-commutative ring]。若 \(\alpha\) 为奇次数类，只能立即推出 \(2\alpha^2=0\)，不能在整数系数下直接除以 \(2\)。当系数域的特征不为 \(2\) 时，才可据此断言所有奇次数类的平方为零。特征 \(2\) 时符号消失，奇次数平方则可能非零；下一节的非定向曲面会给出实例。

\subsection{相对杯积及配对}
\label{g2:subsec:0404:04}

相对上链 \(\varphi\in C^p(X,A;R)\) 可以看成在 \(A\) 中每条链上都取零的 \(X\) 上链。若另一个上链 \(\psi\) 在 \(B\) 上为零，则它们的杯积在完全落入 \(A\) 的单形上为零，也在完全落入 \(B\) 的单形上为零。不过，一条落在 \(A\cup B\) 的单形可能同时穿过两部分。这是相对杯积必须处理的额外问题。

\begin{proposition}[相对杯积]\label{g2:prop:relative-cup-product}
	设 \(X\) 为拓扑空间，\(R\) 为含幺交换环，\(A,B\subset X\) 满足
	\begin{equation}\label{g2:eq:relative-cup-excision-condition}
		A\cup B
		=\operatorname{int}_{A\cup B}A\ \cup\ \operatorname{int}_{A\cup B}B.
	\end{equation}
	则有自然的双线性乘积
	\[
		H^p(X,A;R)\otimes_R H^q(X,B;R)
		\longrightarrow H^{p+q}(X,A\cup B;R).
	\]
	它与绝对杯积及相对群到绝对群的映射相容，并满足分次交换律。特别地，\(A,B\) 都开时满足此条件。取 \(A=B\) 或其中之一为空，总可得到相对上同调的乘法以及绝对上同调对它的作用。
\end{proposition}
\begin{proof}
	先用整数链，令
	\[
		S_*=C_*(A;\mathbb Z)+C_*(B;\mathbb Z)
		\subset C_*(A\cup B;\mathbb Z).
	\]
	由\eqref{g2:eq:relative-cup-excision-condition}和\cref{g2:thm:small-chain-isomorphism}，包含 \(S_*\to C_*(A\cup B;\mathbb Z)\) 诱导同调同构。商复形
	\[
		Q_*=C_*(A\cup B;\mathbb Z)/S_*
	\]
	因此无同调。每个 \(Q_n\) 仍为自由阿贝尔群，以不完全落入 \(A\) 或 \(B\) 的单形为基；其圈群和边界群也自由。\cref{g2:lem:free-complex-splitting}在同调全为零时给出的两项分解，说明 \(Q_*\) 是可缩复形。于是张量 \(R\) 后仍可缩，再取 \(\operatorname{Hom}_R(-,R)\) 也可缩。

	自然满射
	\[
		q:C_*(X;R)/(C_*(A;R)+C_*(B;R))
		\longrightarrow C_*(X,A\cup B;R)
	\]
	的核就是 \(Q_*\otimes_{\mathbb Z}R\)。两侧链模都有由单形得到的自由基，故对应短正合列逐次分裂；对其取 \(\operatorname{Hom}_R(-,R)\) 后仍正合，核复形的可缩性说明 \(q^*\) 诱导上同调同构。

	现在取相对上圈 \(\varphi,\psi\)。前后面公式说明 \(\varphi\smile\psi\) 在 \(C_*(A)+C_*(B)\) 上取零，故给出 \(q\) 的源上的上圈类；经上述 \(q^*\) 的逆同构，得到 \(H^{p+q}(X,A\cup B;R)\) 中的类。更换代表时，\cref{g2:prop:cup-leibniz}中的上边界公式仍在这个商复形上成立，所以乘积良定义。连续映射保持相应子空间时，杯积和 \(q\) 都与之交换；同构的逆也自然，因此得到相对自然性。

	\cref{g2:prop:graded-commutativity}构造的 \(K\sigma\) 的两个因子都留在 \(\sigma\) 的像内。若 \(\sigma\) 落在 \(A\)，\(\varphi\) 使其取值为零；若落在 \(B\)，\(\psi\) 使其取值为零。因此 \((\varphi\otimes\psi)K\) 也在 \(S_*\) 上为零，原来的交换同伦下降到同一个商复形，证明相对分次交换律。与绝对杯积的相容性由各商映射与前后面公式直接得到。
\end{proof}

当 \(A,B,C\) 都为开集时，可以继续相乘并得到 \(H^*(X,A\cup B\cup C;R)\) 中的结合律：三重上链乘积在 \(C_*(A)+C_*(B)+C_*(C)\) 上取零，小链比较与刚才相同，而两种结合方式在这个商上已经相等。这个说法也适用于任意有限个开子空间。不满足小链比较条件的一般子空间对，不能仅凭前后面公式就把值域写成 \(H^*(X,A\cup B;R)\)。

\begin{definition}[取值配对与上同调外积]\label{g2:def:kronecker-cross-cohomology}
	设 \(R\) 为含幺交换环。相对上圈 \(\varphi\) 与同次数相对圈 \(z\) 的\term[quzhi peidui]{取值配对}[evaluation pairing]定义为
	\[
		\langle[\varphi],[z]\rangle=\varphi(z)\in R.
	\]
	对 \(\alpha\in H^p(X;R)\)、\(\beta\in H^q(Y;R)\)，定义\term[shangtongdiao waiji]{上同调外积}[cohomology cross product]为
	\[
		\alpha\times\beta
		=\operatorname{pr}_X^*\alpha\smile\operatorname{pr}_Y^*\beta
		\in H^{p+q}(X\times Y;R).
	\]
	若 \(A\subset X\)、\(B\subset Y\) 为开子空间，同一公式利用相对杯积给出
	\[
		H^p(X,A;R)\otimes_R H^q(Y,B;R)
		\longrightarrow
		H^{p+q}(X\times Y,A\times Y\cup X\times B;R).
	\]
\end{definition}

配对与代表无关：上边界在圈上为零，上圈在边界上为零；相对情形中，剩余链落在子空间，仍被相对上链消去。配对不总能看见所有上同调类。例如 \(H^2(\mathbb{RP}^2;\mathbb Z)=\mathbb Z/2\) 而 \(H_2(\mathbb{RP}^2;\mathbb Z)=0\)，所以这次配对恒为零。万有系数定理中的 Ext 项正是其中的一种不可见信息。\symidx{\(\langle\alpha,z\rangle\)：上同调类与同调类的取值配对}

\begin{proposition}[外积的取值]\label{g2:prop:cross-cup-evaluation}
	设 \(R\) 为含幺交换环，\(\alpha\in H^p(X;R)\)、\(\beta\in H^q(Y;R)\)、\(a\in H_p(X;R)\)、\(b\in H_q(Y;R)\)。则
	\[
		\langle\alpha\times\beta,a\times b\rangle
		=\langle\alpha,a\rangle\langle\beta,b\rangle.
	\]
	若因子配对的次数不匹配而总次数相等，则取值为零。对开子空间 \(A\subset X\)、\(B\subset Y\)，同一结论适用于相对外积和相对配对。另有
	\[
		\Delta_X^*(\alpha\times\beta)=\alpha\smile\beta,
	\]
	以及外积的乘法规则
	\[
		(\alpha\times\beta)\smile(\alpha'\times\beta')
		=(-1)^{|\beta||\alpha'|}
			(\alpha\smile\alpha')\times(\beta\smile\beta').
	\]
\end{proposition}
\begin{proof}
	取代表上圈 \(\varphi,\psi\)。上同调外积的上链代表恰为 \((\varphi\otimes\psi)\operatorname{AW}\)，这可直接对乘积中的单形应用前后面公式。于是
	\[
		\langle\alpha\times\beta,a\times b\rangle
		=(\varphi\otimes\psi)\operatorname{AW}\operatorname{sh}(z\otimes w).
	\]
	由\cref{g2:thm:eilenberg-zilber}，\(\operatorname{AW}\operatorname{sh}\) 与恒等链同伦；\(z\otimes w\) 是圈，\(\varphi\otimes\psi\) 是上圈，所以右侧等于 \(\varphi(z)\psi(w)\)。双次数不匹配时，泛函的定义使它为零。

	相对情形还要核对模型。令 \(Z=A\times Y\cup X\times B\)。两条链映射及其自然链同伦分别保持来自 \(A\times Y\)、\(X\times B\) 的链，因而下降为
	\[
		\frac{C_*(X\times Y;R)}{C_*(A\times Y;R)+C_*(X\times B;R)}
		\ \simeq_{\mathrm{ch}}\
		C_*(X,A;R)\otimes_R C_*(Y,B;R).
	\]
	右侧的商描述由各奇异单形基逐项得到。由于 \(A,B\) 开，左侧分母是 \(Z\) 的开覆盖小链复形。\cref{g2:prop:relative-cup-product}证明中的可缩商论证，给出左侧与 \(C_*(X\times Y,Z;R)\) 的自然同调、上同调比较同构；通过它定义相对同调外积。杯积定义的相对上同调外积在小链商上仍由 \((\varphi\otimes\psi)\operatorname{AW}\) 表示，因此完全相同的链同伦取值证明适用。

	对角公式来自 \(\operatorname{pr}_X\Delta_X=\operatorname{pr}_Y\Delta_X=\operatorname{id}_X\) 及自然性。最后，把四个拉回类按 \(\alpha,\beta,\alpha',\beta'\) 的顺序相乘，只需把 \(\beta\) 与 \(\alpha'\) 交换一次；\cref{g2:prop:graded-commutativity}给出系数 \((-1)^{|\beta||\alpha'|}\)，其余由结合律完成。
\end{proof}

这个配对公式还给出常用的上同调 Künneth 形式。若 \(k\) 为域，且对每个 \(p\)，\(H_p(X;k)\) 有限维，则
\[
	\bigoplus_{p+q=n}H^p(X;k)\otimes_k H^q(Y;k)
	\xrightarrow{\ \times\ }H^n(X\times Y;k)
\]
为自然同构。理由是：域上的万有系数同构把上同调识别为同调的线性对偶，\cref{g2:thm:kunneth}计算乘积同调，而有限维条件保证
\[
	\operatorname{Hom}_k(U\otimes V,k)
	\cong\operatorname{Hom}_k(U,k)\otimes\operatorname{Hom}_k(V,k)
	\quad(\dim U<\infty).
\]
选 \(U\) 的有限基及对偶基，就可把任意双线性泛函写成这类张量的有限和；这也证明所用的代数事实。外积取值公式确认所得同构正是上面显示的映射。若 \(R\) 为主理想整环，两个空间各次同调都为有限秩自由 \(R\) 模，同一论证也适用；Ext、Tor 项在此都为零。其分次环的乘法必须采用命题中的符号规则。

\begin{proposition}[局部生成元的乘积]\label{g2:prop:local-cross-normalization}
	设 \(p,q\ge1\)，在 \(\mathbb R^p,\mathbb R^q\) 上取标准坐标定向，在 \(\mathbb R^{p+q}=\mathbb R^p\times\mathbb R^q\) 上先列第一因子坐标。令
	\[
		u_p\in H^p(\mathbb R^p,\mathbb R^p\setminus\{0\};\mathbb Z),
		\qquad
		u_q\in H^q(\mathbb R^q,\mathbb R^q\setminus\{0\};\mathbb Z)
	\]
	分别在正向局部同调生成元上取值 \(1\)。则相对外积 \(u_p\times u_q\) 是
	\[
		H^{p+q}(\mathbb R^{p+q},\mathbb R^{p+q}\setminus\{0\};\mathbb Z)
	\]
	的正向生成元。这个结论可迭代到任意有限个因子，也适用于原点的球形或立方体邻域。
\end{proposition}
\begin{proof}
	\cref{g2:thm:manifold-local-homology}把这两个局部同调群分别识别为次数 \(p,q\) 的 \(\mathbb Z\)，其他次数为零。万有系数定理因而使相应上同调与同调的对偶一致，故在正向生成元上取值 \(1\) 唯一确定 \(u_p,u_q\)。

	各取一个以原点为内点的小闭立方体，按标准坐标定向三角剖分。把各最高维单形带定向求和，得到链 \(c_p,c_q\)。在“立方体模边界”的相对胞腔链中，这个和代表正向最高维胞腔的生成元；径向投影把穿孔立方体形变收缩到边界，所以由空间对的长正合列和局部切除，它也代表正向局部类。其 shuffle 链是两个剖分的乘积剖分：每个单形乘积按格路径切分，\eqref{g2:eq:shuffle-map}中的符号使所有块带标准乘积定向，内部面在相邻块之间抵消。故 \(\operatorname{sh}(c_p\otimes c_q)\) 正是乘积立方体的正向基本链，代表局部生成元 \(c_{p+q}\)。特别地，这里给出的是正号，而不只是一个未定号的生成元。

	两个穿孔子空间都开，且
	\[
		(\mathbb R^p\setminus\{0\})\times\mathbb R^q
		\ \cup\ \mathbb R^p\times(\mathbb R^q\setminus\{0\})
		=\mathbb R^{p+q}\setminus\{0\}.
	\]
	因此可用\cref{g2:prop:cross-cup-evaluation}的相对形式，得到
	\[
		\langle u_p\times u_q,c_{p+q}\rangle
		=\langle u_p,c_p\rangle\langle u_q,c_q\rangle=1.
	\]
	这证明了归一化。反复应用即得多因子结论；限制到包含原点的球或立方体邻域时，切除给出同一局部群，正向基本链的选取不变。
\end{proof}

例如在 \(\mathbb C^n\) 中采用 \((x_1,y_1,\ldots,x_n,y_n)\) 的实坐标次序，每个坐标复平面的正向局部二次类相乘，给出原点的正向 \(2n\) 次类。相对杯积因此能够把若干坐标超平面的局部条件集中到它们的交点；下一节计算复射影空间的乘法时，正是用这个可直接核验的局部事实来确定系数。

% !TeX root = ../../Book2.tex
% 纲要标题：### 4.5 上同调环的计算与比较
% 纲要位置：几何学计划纲要.md，第 1183 行。
% 正文状态：初稿；尚未正式审核。

% 纲要内容开始
% 
% * 环面和闭曲面的上同调环
% * 复射影空间的生成元和关系
% * 相同加性同调、不同乘法结构的模型
% 
% ---
% 
% 纲要内容结束

\section{上同调环的计算与比较}
\label{g2:sec:0405}

上一章已经算出环面、闭曲面和射影空间的同调群。现在的问题是：两个低次上同调类能否相乘成为高次的非零类？加性群只告诉我们每个次数有哪些元素，杯积还记录这些元素怎样共同作用于空间。下面先从坐标投影计算环面，再用连通和的压缩映射处理闭曲面；复射影空间的关键则是把若干个相对上同调类集中到同一个坐标邻域中。本节写 \(\mathbb F_2=\mathbb Z/2\mathbb Z\)，所有环同构均保留次数和单位元。

\subsection{环面和闭曲面的上同调环}
\label{g2:subsec:0405:01}
取圆周的一次整数上同调生成元 \(\xi\)，使它在正向绕行的圈上取值 \(1\)。在环面 \(T^m=(S^1)^m\) 上，第 \(i\) 个坐标投影 \(p_i\) 给出类 \(x_i=p_i^*\xi\)。这些类分别来自各个圆周因子，它们的乘积可用前两节的外积及其取值公式计算：各因子的一次生成元相乘，检测相应坐标子环面的基本圈。

\begin{proposition}[环面的上同调环]\label{g2:prop:torus-cohomology-ring}
	设 \(m\geq1\)。整数系数下，\(T^m=(S^1)^m\) 的上同调环为
	\[
		H^*(T^m;\mathbb Z)\cong\Lambda_{\mathbb Z}(x_1,\ldots,x_m),\qquad |x_i|=1.
	\]
	右端是由次数为一的生成元及关系 \(x_i^2=0\)、\(x_ix_j=-x_jx_i\) 构成的\term[waidaishu]{外代数}[exterior algebra]；每个次数 \(k\) 的加性基由 \(x_{i_1}\cdots x_{i_k}\)、\(i_1<\cdots<i_k\) 给出。对任意交换含幺环 \(R\)，相同陈述成立，只需将 \(\mathbb Z\) 换成 \(R\)。
\end{proposition}
\begin{proof}
	先取整数系数。圆周的整数同调仅在次数 \(0,1\) 为 \(\mathbb Z\)。由\cref{g2:thm:eilenberg-zilber,g2:thm:kunneth}，反复对圆周作乘积时，外积给出的类在每个次数构成基；这里所有参与的群均为自由群，没有额外的 \(\operatorname{Tor}\) 项。\cref{g2:prop:cross-cup-evaluation}将相应上同调外积识别为上述坐标类的杯积，并证明它们与对应乘积圈互为对偶基。

	因为 \(\xi^2\in H^2(S^1;\mathbb Z)=0\)，自然性给出 \(x_i^2=p_i^*(\xi^2)=0\)。对两个不同的坐标，\cref{g2:prop:graded-commutativity}给出 \(x_ix_j=-x_jx_i\)。这些关系把每个单项式化为一个严格递增指标的单项式，或化为零；而外积已经证明所有递增单项式线性无关，所以没有其他关系。

	对一般 \(R\)，把它先视为阿贝尔群。由于环面的各次整数同调有限生成且自由，\cref{g2:thm:uct-cohomology}说明整数上同调基的系数变换给出 \(H^k(T^m;R)\) 的 \(R\) 基。系数变换与杯积相容，所以刚才的生成元及全部乘法关系随之成立。在特征 \(2\) 时，负号虽然消失，\(x_i^2=0\) 仍是独立的关系，不能只从交换性中推出。
\end{proof}

对于 \(T^2\)，因此有 \(x^2=y^2=0\)，而 \(xy\) 是二次生成元。两个一次类分别测量两个方向的绕行，它们的乘积能够检测整个二维胞腔。若改成两个圆周与一个球面的楔和，虽然各次加性群一样，两个一次类却不能这样相乘；本节末尾会明确比较这两个空间。

\ProseParagraphBreak

为了把环面的计算搬到多柄曲面上，先说明楔和与连通和之间所用的映射。若 \(X_1,\ldots,X_s\) 是连通 CW 复形，各以一个零胞腔为基点，则有限楔和 \(W=\bigvee_iX_i\) 的约化胞腔链分解为各因子的直和。经胞腔比较及自然的万有系数正合列，限制映射给出
\begin{equation}\label{g2:eq:wedge-positive-cohomology}
	\widetilde H^k(W;R)\xrightarrow{\ \cong\ }\bigoplus_i\widetilde H^k(X_i;R),\qquad k>0.
\end{equation}
这里的乘法按分量进行：来自不同因子的两个正次类，其乘积限制到每个因子上都为零，由该同构的单射性，乘积本身就是零。零次则只有一个共同的单位元，而不是每个因子各保留一个单位元。

\begin{lemma}[连通和的压缩映射]\label{g2:lem:surface-pinch-cohomology}
	设 \(M=M_1\mathbin{\#}\cdots\mathbin{\#}M_s\) 是 \(s\geq1\) 个连通闭曲面模型的连通和。若各 \(M_i\) 可定向，给它们选择与连通和相容的定向并取 \(R=\mathbb Z\)；若不要求可定向，则取 \(R=\mathbb F_2\)。存在连续映射 \(q:M\to\bigvee_iM_i\)，使一次拉回
	\[
		q^*:H^1\Bigl(\bigvee_iM_i;R\Bigr)\longrightarrow H^1(M;R)
	\]
	为同构。若 \(\omega_i\in H^2(M_i;R)\) 是对应定向的生成元，模 \(2\) 时取唯一非零类，则各 \(\omega_i\) 在楔和中的延拓均被 \(q^*\) 送到同一个二次生成元 \(\omega\in H^2(M;R)\)。
\end{lemma}
\begin{proof}
	先考虑两个曲面的连通和。将粘合用的分离圆周压成一点，左右两块各自成为 \((M_i\setminus\operatorname{int}D^2)/\partial D^2\)。这与 \(M_i\) 同胚：边界的圆环领圈在内边界压缩后变成一个圆盘，外边界上的粘合保持不变。于是得到到 \(M_1\vee M_2\) 的映射。多次作此操作，得到各因子由若干基点连成树的空间；再压掉连接这些基点的树，可把它们归并为一个楔点。等价地，可在原曲面的各连接圆周之间选互不交叉的连接弧，一并压缩。这个构造在每个因子内保留一组标准边环路。

	由\eqref{g2:eq:surface-attaching-words}，可定向曲面的整数一维同调以各柄的两条边为基；模 \(2\) 时，无论可否定向，各条标准边也独立生成一维同调，因为二维胞腔的边界系数均为偶数。\(q_*\) 将这些边送到相应因子的同名边，故在一次同调上是同构。\cref{g2:thm:uct-cohomology}在一次给出自然同构 \(H^1(-;R)\cong\operatorname{Hom}_{\mathbb Z}(H_1(-;\mathbb Z),R)\)，因为零次同调是自由群；这就证明了关于 \(q^*\) 的断言。

	还要确定二维生成元的倍数。将多边形模型细分为与连接圆周及连接弧相容的有限三角剖分；可先切分各边和各多边形，再作重心细分，使粘合后的三角形有一致的顶点次序。可定向时按曲面定向给各三角形取符号，模 \(2\) 时全部取系数 \(1\)。这些三角形之和 \(z\) 是奇异二圈：内部边出现两次并抵消，原多边形的成对边也按所取系数抵消。相对化到原二维胞腔时，这些小三角片各出现一次并铺满特征圆盘，所以其胞腔系数为 \(1\)；胞腔比较定理因而将 \([z]\) 识别为二维生成元。把 \(q\) 再与到第 \(i\) 个因子的收缩复合，记为 \(q_i:M\to M_i\)。在该因子内选一个避开压缩集合的三角形内点 \(y\)，其原像只有一个点，且 \(q_i\) 在此处保持定向地为局部同胚；模 \(2\) 时无须选择符号。

	映到局部同调 \(H_2(M_i,M_i\setminus\{y\};R)\) 后，\((q_i)_*[z]\) 的系数因此为 \(1\)。同一个三角形和构造说明 \(M_i\) 的二维生成元也映到局部生成元：切除到该三角形内部，只剩包含 \(y\) 的那一片，边界定向与圆盘相对生成元一致。由于两边的二次群均为秩一的自由 \(R\) 模，\((q_i)_*[z]\) 恰为所选生成元。最后用取值配对的自然性得到 \(\langle q_i^*\omega_i,[z]\rangle=1\)。可定向情形下一维整数同调自由，模 \(2\) 情形下系数是域，因此相应的上同调取值同构保证 \(q_i^*\omega_i=\omega\)。
\end{proof}

\begin{theorem}[可定向曲面的上同调环]\label{g2:thm:orientable-surface-ring}
	设 \(\Sigma_g\) 是 \(g\geq1\) 个环面的连通和，并选定定向。存在一次类 \(a_1,b_1,\ldots,a_g,b_g\) 和二次类 \(\omega\)，使
	\[
		H^0(\Sigma_g;\mathbb Z)=\mathbb Z,
		\quad H^1(\Sigma_g;\mathbb Z)=\bigoplus_{i=1}^g(\mathbb Za_i\oplus\mathbb Zb_i),
		\quad H^2(\Sigma_g;\mathbb Z)=\mathbb Z\omega.
	\]
	所有更高次群为零，正次生成元的乘法由
	\begin{equation}\label{g2:eq:orientable-surface-products}
		a_i a_j=b_i b_j=0,\qquad a_i b_j=\delta_{ij}\omega,\qquad b_j a_i=-\delta_{ij}\omega
	\end{equation}
	完全决定。\(\Sigma_0=S^2\) 的环为 \(\mathbb Z[\omega]/(\omega^2)\)，其中 \(|\omega|=2\)。
\end{theorem}
\begin{proof}
	取\cref{g2:lem:surface-pinch-cohomology}的映射 \(q:\Sigma_g\to\bigvee_iT^2\)。在每个环面上选择\cref{g2:prop:torus-cohomology-ring}的一次基 \(x_i,y_i\)，使 \(x_i y_i=\omega_i\)，并通过收缩其他因子把这些类延拓到楔和。令 \(a_i=q^*x_i\)、\(b_i=q^*y_i\)。同因子中的平方为零，异因子之间的全部正次乘积为零；而 \(q^*\omega_i=\omega\)。杯积自然性于是给出全部关系。加性群已由\cref{g2:prop:surface-integer-homology}和万有系数定理确定，所以这些关系没有遗漏其他乘法。\(S^2\) 的结论直接来自它的上同调只在次数 \(0,2\) 非零。
\end{proof}

不可定向曲面宜先用模 \(2\) 系数观察乘法。这样即使一条边沿同一个方向出现两次，它在边界中仍会抵消。与可定向曲面的主要区别不是交换律中的负号，而是某个一次类的平方可以非零。我们先把这个现象在 \(\mathbb RP^2\) 上算清楚。

\begin{lemma}[射影平面的一次平方]\label{g2:lem:projective-plane-square}
	设 \(t\) 是 \(H^1(\mathbb RP^2;\mathbb F_2)\) 的非零类，则 \(t^2\) 是 \(H^2(\mathbb RP^2;\mathbb F_2)\) 的非零类。因此
	\[
		H^*(\mathbb RP^2;\mathbb F_2)\cong\mathbb F_2[t]/(t^3),\qquad |t|=1.
	\]
\end{lemma}
\begin{proof}
	考虑系数短正合列
	\[
		0\longrightarrow\mathbb F_2\xrightarrow{\ \iota\ }\mathbb Z/4\xrightarrow{\ \rho\ }\mathbb F_2\longrightarrow0,
		\qquad\iota(1)=2,
	\]
	其中 \(\rho\) 为模 \(2\) 映射。奇异链群是自由阿贝尔群，所以对每个奇异单形逐一提升上链的值，得到上链复形的短正合列及连接同态 \(\beta:H^1(X;\mathbb F_2)\to H^2(X;\mathbb F_2)\)。它的具体算法是：把模 \(2\) 闭上链 \(c\) 的每个值提升为 \(0\) 或 \(1\)，得到模 \(4\) 上链 \(\widetilde c\)；其上边界取值于 \(\{0,2\}\)，除以 \(2\) 再模 \(2\)，所得闭上链代表 \(\beta[c]\)。

	在一个奇异二单形 \([v_0v_1v_2]\) 上，记 \(c[v_0v_1]=r\)、\(c[v_1v_2]=s\)，其中 \(r,s\in\{0,1\}\)。闭性条件给出 \(c[v_0v_2]\equiv r+s\pmod2\)，因此
	\[
		\frac{(\delta\widetilde c)[v_0v_1v_2]}2
		=\frac{s-\bigl((r+s)\bmod2\bigr)+r}{2}
		\equiv rs\pmod2.
	\]
	右端正是\cref{g2:def:cup-product}的 \((c\smile c)[v_0v_1v_2]\)。于是对任意空间，\(\beta[c]=[c]^2\) 在一次都成立。

	对 \(X=\mathbb RP^2\)，\cref{g2:prop:projective-integer-homology}给出 \(H_1(X;\mathbb Z)=\mathbb Z/2\)。一次万有系数同构及其自然性说明 \(H^1(X;\mathbb Z/4)\to H^1(X;\mathbb F_2)\) 对应 \(\operatorname{Hom}(\mathbb Z/2,\mathbb Z/4)\to\operatorname{Hom}(\mathbb Z/2,\mathbb F_2)\)。前一群的同态只能把生成元送到 \(0\) 或 \(2\)，经 \(\rho\) 都变成零。系数长正合列因而说明 \(\beta\) 在一次单射。由\cref{g2:cor:cellular-coefficients}，模 \(2\) 胞腔微分全为零；域系数的上同调取值同构因而说明 \(H^1(X;\mathbb F_2)\) 与 \(H^2(X;\mathbb F_2)\) 都是一维；于是 \(t^2=\beta(t)\ne0\)。高于二次的群为零，得到所述环。
\end{proof}

\begin{proposition}[不可定向曲面的模二上同调环]\label{g2:prop:nonorientable-surface-ring}
	设 \(N_h\) 是 \(h\geq1\) 个实射影平面的连通和。存在一次基 \(c_1,\ldots,c_h\) 及二次生成元 \(\omega\)，使
	\[
		H^0(N_h;\mathbb F_2)=\mathbb F_2,\qquad
		H^1(N_h;\mathbb F_2)=\mathbb F_2^h,\qquad
		H^2(N_h;\mathbb F_2)=\mathbb F_2\omega,
	\]
	更高次群为零，而乘法为
	\begin{equation}\label{g2:eq:nonorientable-surface-products}
		c_i c_j=0\quad(i\ne j),\qquad c_i^2=\omega\quad(1\leq i\leq h).
	\end{equation}
\end{proposition}
\begin{proof}
	\cref{g2:lem:surface-pinch-cohomology}给出 \(q:N_h\to\bigvee_i\mathbb RP^2\)，且一次拉回同构，二次拉回将各因子的非零类送到共同的 \(\omega\)。把各因子的一次类记为 \(t_i\)，令 \(c_i=q^*t_i\)。由\cref{g2:lem:projective-plane-square}，\(c_i^2=q^*(t_i^2)=\omega\)；异因子的乘积由\eqref{g2:eq:wedge-positive-cohomology}为零。这给出了全部一次乘法，而更高正次乘积因次数超过 \(2\) 为零。
\end{proof}

这里的系数必须明确。整数系数下，\(H^2(N_h;\mathbb Z)=\mathbb Z/2\) 来自万有系数正合列中的 \(\operatorname{Ext}\) 项；模 \(2\) 系数下，顶次类又能通过三角形和来评价。它们不是可以不加说明地互换的同一个计算。特别地，模 \(2\) 交换律并不迫使奇次平方为零，\(c_i^2=\omega\) 正说明了这一点。

\subsection{复射影空间的生成元和关系}
\label{g2:subsec:0405:02}
\cref{g2:prop:projective-integer-homology}已经给出 \(\mathbb CP^n\) 的整数同调：每个偶数次数 \(0,2,\ldots,2n\) 各有一个 \(\mathbb Z\)，其余为零。由于这些群自由，万有系数定理给出完全相同的上同调加性群。但是，仅凭这个结果还不能断言二次生成元 \(u\) 的平方是四次生成元：它可能是零，也可能是该生成元的某个非单位整数倍。下面的证明正是为了确定这个整数倍。

\begin{theorem}[复射影空间的上同调环]\label{g2:thm:complex-projective-ring}
	对整数 \(n\geq1\)，若 \(u\) 是 \(H^2(\mathbb CP^n;\mathbb Z)\) 的任一生成元，则
	\[
		H^*(\mathbb CP^n;\mathbb Z)\cong\mathbb Z[u]/(u^{n+1}),\qquad |u|=2.
	\]
	标准包含 \(\mathbb CP^k\hookrightarrow\mathbb CP^n\) 在 \(0\leq k\leq n\) 时将该环限制为截断到次数 \(2k\) 的环。赋予标准 CW 弱拓扑的无限复射影空间满足
	\[
		H^*(\mathbb CP^\infty;\mathbb Z)\cong\mathbb Z[u],\qquad |u|=2.
	\]
\end{theorem}
\begin{proof}
	先证有限维情形。记 \(X=\mathbb CP^n\)，在齐次坐标 \([z_0:\cdots:z_n]\) 中取点 \(p=[1:0:\cdots:0]\) 及 \(n\) 个超平面
	\[
		H_j=\{z_j=0\},\qquad A_j=X\setminus H_j=\{z_j\ne0\},\qquad 1\leq j\leq n.
	\]
	每个 \(A_j\) 都是仿射坐标图 \(\mathbb C^n\)，因而可缩。\cref{g2:prop:relative-cohomology-exactness}的空间对长正合列说明忘去相对条件的映射 \(H^2(X,A_j;\mathbb Z)\to H^2(X;\mathbb Z)\) 是同构，因为它两侧的 \(H^1(A_j;\mathbb Z)\)、\(H^2(A_j;\mathbb Z)\) 都为零。于是 \(u\) 有唯一相对提升 \(\widetilde u_j\)。这些提升可以理解为：计算 \(u\) 时，把超平面 \(H_j\) 之外的部分作为可以忽略的区域。

	相对杯积的目标由被忽略区域的并集决定。所有 \(A_j\) 都开，且
	\[
		A_1\cup\cdots\cup A_n=X\setminus\{p\}.
	\]
	由\cref{g2:prop:relative-cup-product}，迭代相对杯积得到
	\[
		v=\widetilde u_1\smile\cdots\smile\widetilde u_n
		\in H^{2n}(X,X\setminus\{p\};\mathbb Z),
	\]
	它在绝对上同调中的像为 \(u^n\)。证明现在分成两件事：\(v\) 是局部群的生成元，而且这个局部群到绝对顶次上同调的映射是同构。

	取含 \(p\) 的仿射图 \(U=\{z_0\ne0\}\)，以 \(w_j=z_j/z_0\) 识别为 \(\mathbb C^n\)。在这里，\(A_j\cap U=\{w_j\ne0\}\)。要知道 \(\widetilde u_j\) 限制后的整数倍数，考虑坐标直线
	\[
		L_j=\{[z_0:0:\cdots:z_j:\cdots:0]\}\cong\mathbb CP^1.
	\]
	标准胞腔包含在二次同调上为同构，故 \(H^2(X;\mathbb Z)\to H^2(L_j;\mathbb Z)\) 也是同构。对于非第一条坐标直线，交换坐标给出同样的结论；这种交换是 \(X\) 的同胚，并把第一条坐标直线送到 \(L_j\)。又因 \(A_j\cap L_j=L_j\setminus\{p\}\cong\mathbb C\)，下图两条水平箭头和右侧箭头均为同构：
	\[
		\begin{tikzcd}[column sep=large]
			H^2(X,A_j;\mathbb Z) \arrow[r] \arrow[d] & H^2(X;\mathbb Z) \arrow[d] \\
			H^2(L_j,L_j\setminus\{p\};\mathbb Z) \arrow[r] & H^2(L_j;\mathbb Z).
		\end{tikzcd}
	\]
	因此左侧箭头也为同构，\(\widetilde u_j\) 限制到这条直线的局部类是生成元。再切除到 \(L_j\cap U\)，得到 \(H^2(\mathbb C,\mathbb C\setminus\{0\};\mathbb Z)\) 的生成元。另一方面，空间对 \((U,A_j\cap U)\) 沿其他 \(n-1\) 个坐标同时收缩到这条坐标轴，故它的相对二次上同调也由该限制完全确定。若 \(\eta_j\) 表示 \(\mathbb C\) 在原点处的一个相对生成元，\(\pi_j:U\to\mathbb C\) 为坐标投影，则 \(\widetilde u_j|_U=\pm\pi_j^*\eta_j\)。

	于是 \(v|_U\) 就是 \(n\) 个平面局部生成元的外积，至多相差符号。\cref{g2:prop:local-cross-normalization}说明，该外积与乘积圆盘的相对基本链配对为 \(\pm1\)，所以它生成
	\[
		H^{2n}(\mathbb C^n,\mathbb C^n\setminus\{0\};\mathbb Z)\cong\mathbb Z.
	\]
	这一步也可直接把每个平面换成原点周围的小闭圆盘：各相对二次类在圆盘上取值 \(\pm1\)，乘积链的评价是这些数的乘积，而乘积圆盘的边界恰对应“至少一个坐标落在边界上”。\cref{g2:prop:relative-cohomology-exactness}的切除结论把这个局部群与 \(H^{2n}(X,X\setminus\{p\};\mathbb Z)\) 识别，故 \(v\) 是后者的生成元。

	最后，\(X\setminus\{p\}\) 变形收缩到超平面 \(\{z_0=0\}\cong\mathbb CP^{n-1}\)。显式收缩为
	\[
		[z_0:z_1:\cdots:z_n]\longmapsto[(1-t)z_0:z_1:\cdots:z_n],\qquad 0\leq t\leq1.
	\]
	在所用定义域中，\((z_1,\ldots,z_n)\ne0\)，所以此公式始终有定义并固定目标超平面。因而补空间在次数 \(2n-1,2n\) 的上同调均为零，空间对长正合列给出
	\[
		H^{2n}(X,X\setminus\{p\};\mathbb Z)\xrightarrow{\ \cong\ }H^{2n}(X;\mathbb Z).
	\]
	\(v\) 的像 \(u^n\) 因此是顶次生成元。

	对每个 \(1\leq k\leq n\)，标准包含 \(\mathbb CP^k\hookrightarrow\mathbb CP^n\) 在二次和 \(2k\) 次上同调上都是同构；这由相容的偶维胞腔基及万有系数同构得到。将上述顶次计算用于 \(\mathbb CP^k\)，可知 \(u^k\) 生成 \(H^{2k}(\mathbb CP^n;\mathbb Z)\)。这些幂连同单位元已经穷尽全部加性基，且 \(2n\) 次以上没有群，故唯一所需截断关系是 \(u^{n+1}=0\)。

	对于 \(\mathbb CP^\infty\)，每个固定次数的胞腔链和对应同调在足够大的有限阶段已经稳定；相应同调群自由，万有系数同构于是保证 \(H^{2k}(\mathbb CP^\infty;\mathbb Z)\to H^{2k}(\mathbb CP^n;\mathbb Z)\) 在 \(n\geq k\) 时为同构。选择二次生成元 \(u\)，再限制到各有限阶段，就知每个 \(u^k\) 都生成该次群。分次上同调环按各次群的直和定义，一个元素只有有限多个非零齐次分量，因此得到多项式环 \(\mathbb Z[u]\)，而不是形式幂级数环。
\end{proof}

证明中起决定作用的是超平面的共同交点。在 \(\mathbb CP^2\) 中，两个坐标射影直线交于 \(p\)，它们对应的相对二次类在 \(p\) 附近分别检测两个复坐标方向；相乘后就检测一个四维邻域。这个图景解释了 \(u^2\) 为什么有机会非零，而局部外积配对为 \(\pm1\) 则进一步排除了它只是顶次生成元的两倍或三倍。这里的坐标与相对杯积证明可与\cite[第3.2节，Spaces with Polynomial Cohomology]{Hatcher2002}对照。

\subsection{相同加性同调、不同乘法结构的模型}
\label{g2:subsec:0405:03}
前面的计算给出两种区分空间的办法：环面的一次类可以相乘得到顶次类，复射影平面的二次类可以平方得到四次生成元。下面分别与加性同调相同的楔和比较，考察这两种乘积在楔和中为何消失。

\begin{example}[环面与球面圆周的楔和]\label{g2:exm:torus-wedge-ring}
	令 \(X=T^2\)、\(Y=S^1\vee S^1\vee S^2\)。两者的整数同调群在次数 \(0,1,2\) 分别是 \(\mathbb Z,\mathbb Z^2,\mathbb Z\)，其他次数为零。但 \(H^*(X;\mathbb Z)\) 中存在一次类 \(x,y\) 使 \(xy\ne0\)，而 \(H^*(Y;\mathbb Z)\) 中任意两个一次类的乘积均为零。后一断言可由\eqref{g2:eq:wedge-positive-cohomology}逐因子核对：在圆周因子上二次群为零，在球面因子上一次群为零。于是两环不可能同构。因为连续映射诱导保持杯积的拉回，同伦等价必诱导上同调环同构，所以 \(X,Y\) 不同伦等价。
\end{example}

\begin{example}[复射影平面与两个球面的楔和]\label{g2:exm:projective-plane-wedge-ring}
	令 \(X=\mathbb CP^2\)、\(Y=S^2\vee S^4\)。两者的整数同调群都仅在次数 \(0,2,4\) 为 \(\mathbb Z\)。然而 \(X\) 的二次生成元 \(u\) 满足 \(u^2\ne0\)，而 \(Y\) 的每个二次类平方为零：它在 \(S^2\) 上的平方因次数为零，在 \(S^4\) 上的限制本身就为零。故两者不同伦等价。这里没有一维同调参与，差异完全发生在偶数次的乘法中。
\end{example}

系数变化还可能消去加性群的区别，同时保留乘法上的差异。对 \(g\geq1\)，\(\Sigma_g\) 与 \(N_{2g}\) 的模 \(2\) 加性同调都在次数 \(0,1,2\) 为 \(\mathbb F_2,\mathbb F_2^{2g},\mathbb F_2\)。前者任意一次类的平方为零：\eqref{g2:eq:orientable-surface-products}模 \(2\) 后，每个基元平方为零，交叉项成对出现。后者则有 \(c_i^2=\omega\ne0\)，所以两者的模 \(2\) 上同调环不同。这项比较的系数始终为 \(\mathbb F_2\)。


\begin{chaptersummary}
	自由消解把一个模写成自由生成元及其关系，并通过比较链映射和比较同伦保证计算不依赖所选表示。取同态后的一次上同调是 Ext，取张量后的一次同调是 Tor；在主理想整环上，自由模的子模仍自由，所以所需消解可以缩短到两项。

	\ProseParagraphBreak

	自由链复形的圈、边界及链群之间有两条不同的短正合列：链群到低一次边界的列可以分裂，边界到圈的列却可能留下挠同调。把两者结合便得到万有系数定理。上同调的取值映射会漏掉低一次同调的 Ext，系数同调中的旧圈张量则会漏掉低一次同调的 Tor。正合列中的映射自然，选出的直和分解一般不自然。

	\ProseParagraphBreak

	Alexander--Whitney 映射沿前后面切开单形，shuffle 映射把单形乘积按格路径剖分。标准模型中的逐次填圈给出自然链同伦，因而空间乘积可以用张量复形计算。Künneth 公式中的张量项来自圈的外积，Tor 项还使用挠圈整数倍的填充，两者在总次数上相差一。

	\ProseParagraphBreak

	杯积是上同调外积沿对角映射的拉回。前后面公式在上链上给出结合律、单位和自然性，而交换律需要一个链同伦，并留下 \((-1)^{pq}\) 的符号。相对杯积还要检查小链比较条件；局部外积配对使各个坐标方向的生成元相乘后仍为乘积邻域的生成元。

	\ProseParagraphBreak

	环面的环是一次生成元的外代数，可定向曲面的乘法记录成对的柄方向，不可定向曲面的模二一次平方可以非零。复射影空间中，坐标超平面的相对类相乘集中到共同交点，证明二次生成元的每个允许幂都是生成元。这些乘法关系既能区分加性同调相同的空间，也能限制连续映射在上同调上的作用。
\end{chaptersummary}

\BookExercises

\begin{exercise}\label{g2:ex:cyclic-coefficient-maps}
	令 \(u:\mathbb Z/12\mathbb Z\to\mathbb Z/4\mathbb Z\) 把 \(1\) 的类送到 \(1\) 的类，令 \(G=\mathbb Z/6\mathbb Z\)。写出 \(u\) 在两项自由消解上的一个提升，并计算
	\[
		u_*:\operatorname{Tor}^{\mathbb Z}_1(\mathbb Z/12\mathbb Z,G)\longrightarrow\operatorname{Tor}^{\mathbb Z}_1(\mathbb Z/4\mathbb Z,G)
	\]
	及
	\[
		u^*:\operatorname{Ext}_{\mathbb Z}^1(\mathbb Z/4\mathbb Z,G)\longrightarrow\operatorname{Ext}_{\mathbb Z}^1(\mathbb Z/12\mathbb Z,G).
	\]
	说明为何仅知道两端群的同构类型，还不足以给出这两个映射。
\end{exercise}
\begin{solution}
	源、靶的消解分别为 \(\mathbb Z\xrightarrow{12}\mathbb Z\)、\(\mathbb Z\xrightarrow{4}\mathbb Z\)，右侧项为零次。取零次提升 \(f_0=1\)，一次提升 \(f_1=3\)，因为 \(4f_1=f_0\,12\)，且零次诱导的商映射正是 \(u\)。

	取张量 \(G\) 后，Tor 分别是 \(G[12]=G\) 与 \(G[4]=\{0,3\}\)。诱导映射由一次提升给出，故 \(u_*(g)=3g\)。它满射，核为 \(\{0,2,4\}\)，而不是把 \(G[12]\) 不经计算地包含到 \(G[4]\) 中。

	取同态后箭头反向。一次上链映射仍由预复合乘 \(3\) 给出，所以
	\[
		u^*:G/4G\longrightarrow G/12G,\qquad[g]\longmapsto[3g].
	\]
	这里 \(4G=\{0,2,4\}\)、\(12G=0\)，故源为二阶群，靶为 \(G\)，该映射将非零类送到 \(3\)，是单射。两处的乘 \(3\) 来自同一个消解提升，但 Tor 在核之间协变作用，Ext 在余核之间反变作用。只有群的阶或同构类型不能确定箭头方向及具体倍数。
\end{solution}

\begin{exercise}\label{g2:ex:uct-needs-free-chains}
	固定 \(m\ge2\)，令整数链复形 \(C\) 仅有 \(C_0=\mathbb Z/m\mathbb Z\) 非零。分别计算 \(H^1(C;\mathbb Z)\) 与 \(H_1(C;\mathbb Z/m\mathbb Z)\)。证明：如果删去两种万有系数定理中“链群自由”的条件，两条短正合列都会在这个例子中失败。
\end{exercise}
\begin{solution}
	原复形只在零次有 \(H_0(C)=\mathbb Z/m\mathbb Z\)。对 \(\operatorname{Hom}(C,\mathbb Z)\)，零次项也为零，因为有限阶元素只能映到 \(0\)；故 \(H^1(C;\mathbb Z)=0\)。如果仍套用上同调万有系数列的 \(n=1\) 情形，就会要求非零群
	\[
		\operatorname{Ext}_{\mathbb Z}^1(H_0(C),\mathbb Z)=\mathbb Z/m\mathbb Z
	\]
	单射到零群中，矛盾。

	张量 \(G=\mathbb Z/m\mathbb Z\) 后，复形仍只在零次非零，所以 \(H_1(C;G)=0\)。如果套用同调万有系数列的 \(n=1\) 情形，就会要求零群满射到
	\[
		\operatorname{Tor}^{\mathbb Z}_1(H_0(C),G)=G[m]=G\ne0,
	\]
	也矛盾。这里直接对非自由复形作张量或取同态，与先将模替换成自由消解再作运算并不是同一件事；失败发生在定理的输入条件，而不是 Ext 或 Tor 的循环模公式。
\end{solution}

\begin{exercise}\label{g2:ex:coefficient-boundary-matrix}
	令自由整数链复形 \(C\) 的非零项为 \(C_2=\mathbb Z\langle c\rangle\)、\(C_1=\mathbb Z\langle x,y\rangle\)、\(C_0=\mathbb Z\langle v\rangle\)，规定 \(\mathrm{d}_2c=4x+6y\)、\(\mathrm{d}_1=0\)。计算整数同调，并在 \(G=\mathbb Z/6\mathbb Z\) 时计算全部 \(H_n(C;G)\)、\(H^n(C;G)\)。给出代表元或坐标描述，并指出万有系数定理中的每个非零项。
\end{exercise}
\begin{solution}
	向量 \(a=2x+3y\)、\(b=x+y\) 构成 \(C_1\) 的一组整数基，因为坐标矩阵的行列式为 \(2-3=-1\)。在此基下 \(\mathrm{d}_2c=2a\)，故
	\[
		H_0(C)=\mathbb Z,\qquad H_1(C)=\mathbb Z\langle[b]\rangle\oplus(\mathbb Z/2\mathbb Z)\langle[a]\rangle,\qquad H_n(C)=0\quad(n\ge2).
	\]
	在 \(G\) 中，\(\mathrm{d}_2(tc)=2t\,a\)，所以二次圈由 \(3c\) 生成，得到 \(H_2(C;G)\cong\mathbb Z/2\mathbb Z\)。一次同调为
	\[
		H_1(C;G)=(G/2G)\langle[a]\rangle\oplus G\langle[b]\rangle\cong\mathbb Z/2\mathbb Z\oplus\mathbb Z/6\mathbb Z,
	\]
	零次同调为 \(G\)，其余次数为零。UCT 中一次项全部来自 \(H_1(C)\otimes G\)，而二次项全部来自 \(\operatorname{Tor}_1(H_1(C),G)=G[2]\)。

	令 \(a^*,b^*\) 表示这组基的 \(G\) 值坐标函数。一次上链 \(r a^*+s b^*\) 的上边界为 \(2r\) 倍的二次坐标函数，所以它闭当且仅当 \(r\in\{0,3\}\)，而 \(s\) 任意。由于零次上链微分为零，
	\[
		H^0(C;G)=G,\qquad H^1(C;G)=G[2]\langle a^*\rangle\oplus G\langle b^*\rangle,\qquad H^2(C;G)=G/2G.
	\]
	在 \(H^1\) 的第一项中，\(G[2]\langle a^*\rangle\) 指允许的系数 \(r\) 属于 \(G[2]\)，其非零类具体为 \(3a^*\)。其他次数为零。UCT 中 \(H^1\) 全部来自 \(\operatorname{Hom}(H_1(C),G)\)，因为 \(H_0(C)\) 自由；\(H^2\) 全部来自 \(\operatorname{Ext}^1(H_1(C),G)=G/2G\)。低一次同调中的同一个二阶挠因子，因而分别在二次同调和二次上同调中通过核与余核出现。
\end{solution}

\begin{exercise}\label{g2:ex:mod-p-betti-count}
	设 \(C\) 为非负自由整数链复形，\(n\ge0\)，并且 \(H_n(C)\)、\(H_{n-1}(C)\) 都有限生成。将其写为自由部分与有限循环群之直和，记自由秩为 \(r_j\)，记该分解中阶被素数 \(p\) 整除的非零循环因子数为 \(t_j(p)\)，约定 \(r_{-1}=t_{-1}(p)=0\)。证明
	\[
		\dim_{\mathbb F_p}H_n(C;\mathbb F_p)=r_n+t_n(p)+t_{n-1}(p),\qquad
		\dim_{\mathbb Q}H_n(C;\mathbb Q)=r_n.
	\]
	再证明同样的两个维数公式对 \(H^n(C;\mathbb F_p)\)、\(H^n(C;\mathbb Q)\) 成立。最后说明这些维数不能区分 \(p\) 阶与 \(p^2\) 阶的循环挠因子。
\end{exercise}
\begin{solution}
	一个自由因子 \(\mathbb Z\) 张量 \(\mathbb F_p\) 贡献一个 \(\mathbb F_p\)，而其一次 Tor 为零。对循环因子 \(\mathbb Z/a\mathbb Z\)，\cref{g2:prop:cyclic-ext-tor}给出张量为 \(\mathbb F_p/a\mathbb F_p\)、Tor 为 \(\mathbb F_p[a]\)。若 \(p\mid a\)，乘 \(a\) 为零，两项各为一维；若 \(p\nmid a\)，乘 \(a\) 可逆，两项均为零。有限直和逐项计算，再由同调万有系数短正合列两端维数相加，就得到第一式。对于 \(\mathbb Q\)，每个非零整数的乘法可逆，所以所有循环因子的张量与 Tor 均为零，只剩自由秩 \(r_n\)。

	上同调公式中，一个自由因子的 Hom 贡献一维，Ext 为零。循环因子的 Hom 为 \(\mathbb F_p[a]\)，Ext 为 \(\mathbb F_p/a\mathbb F_p\)，所以同样各由是否 \(p\mid a\) 决定。上同调万有系数列给出相同的维数。取有理系数时，有限阶群到 \(\mathbb Q\) 的同态为零，而 \(\mathbb Q/a\mathbb Q=0\)，故仍只剩自由秩。

	以上维数只检测循环因子的阶是否被 \(p\) 整除，不检测被 \(p\) 整除的次数。例如分别以乘 \(p\)、乘 \(p^2\) 为微分的两项自由复形，整数零次同调分别为 \(\mathbb Z/p\mathbb Z\)、\(\mathbb Z/p^2\mathbb Z\)，但张量 \(\mathbb F_p\) 后两者微分都为零，所有模 \(p\) 同调维数相同；有理系数也都不能区分它们。
\end{solution}

\begin{exercise}\label{g2:ex:lens-product-groups}
	设 \(m,n\ge2\)，\(g=\gcd(m,n)\)。计算 \(L(m,1)\times L(n,1)\) 的全部整数同调群，并指出哪些项来自 Künneth 公式中的 Tor。所写的直和分解是否都由外积自然给出？
\end{exercise}
\begin{solution}
	由\cref{g2:prop:lens-cellular-model}，\(L(m,1)\) 的整数同调在次数 \(0,1,3\) 分别为 \(\mathbb Z,\mathbb Z/m,\mathbb Z\)，其余为零；另一个因子把 \(m\) 换成 \(n\)。自由因子不贡献 Tor，两个挠群的张量与 Tor 均为 \(\mathbb Z/g\)。因此
	\[
		H_j(L(m,1)\times L(n,1);\mathbb Z)\cong
		\begin{cases}
			\mathbb Z,&j=0,6,\\
			\mathbb Z/m\oplus\mathbb Z/n,&j=1,4,\\
			\mathbb Z/g,&j=2,\\
			\mathbb Z^2\oplus\mathbb Z/g,&j=3,\\
			0,&j=5\text{ 或 }j>6.
		\end{cases}
	\]
	其中次数二的 \(\mathbb Z/g\) 来自 \(H_1\otimes H_1\)；次数三的 \(\mathbb Z/g\) 才来自 \(\operatorname{Tor}_1(H_1,H_1)\)。次数四的两个挠项来自 \(H_1\otimes H_3\)、\(H_3\otimes H_1\)，不是 Tor。除次数三以外，所有显示的群都由外积自然识别；次数三应保留自然短正合列
	\[
		0\longrightarrow\mathbb Z^2\longrightarrow H_3
		\longrightarrow\mathbb Z/g\longrightarrow0.
	\]
	次数三的挠类不是两个同调类的外积。不过，本题还有一个比一般分裂定理更强的特点：左端自由、右端有限，而正合列已知分裂，因此 \(H_3\) 的挠子群经右箭头同构到 \(\mathbb Z/g\)。其逆给出这条列的唯一分裂，因为两个分裂之差须为 \(\mathbb Z/g\to\mathbb Z^2\) 的同态，只能为零。所以本题的第三次直和可由外积子群与挠子群共同确定。若 \(g=1\)，相应循环群为零，所有次数都由外积给出。
\end{solution}

\begin{exercise}\label{g2:ex:explicit-mixed-tor-cycle}
	令 \(E\) 的基为 \(x\in E_1\)、\(a\in E_0\)，微分 \(\mathrm{d}x=ma\)；令 \(F\) 的基为 \(y\in F_1\)、\(b\in F_0\)，微分 \(\mathrm{d}y=nb\)，其中 \(m,n\ge2\)，其他次数为零。直接计算 \(E\otimes F\) 的同调，写出次数一的生成圈，并计算交换链映射 \(\tau:E\otimes F\to F\otimes E\) 对所选生成圈的作用。
\end{exercise}
\begin{solution}
	置 \(g=\gcd(m,n)\)。次数零以 \(a\otimes b\) 为基，来自次数一的边界系数是 \(m,n\)，故 \(H_0=\mathbb Z/g\)。次数一的链
	\[
		u(x\otimes b)+v(a\otimes y)
	\]
	为圈当且仅当 \(mu+nv=0\)。除以 \(g\) 后 \(m/g,n/g\) 互素，因此所有整数解恰为
	\[
		u=k\,n/g,\qquad v=-k\,m/g,\qquad k\in\mathbb Z.
	\]
	于是圈群由
	\[
		z=\frac ng\,x\otimes b-\frac mg\,a\otimes y
	\]
	生成。而
	\[
		\mathrm{d}(x\otimes y)=m\,a\otimes y-n\,x\otimes b=-gz,
	\]
	故 \(H_1=\mathbb Z/g\)。这个边界同态从次数二的 \(\mathbb Z\) 单射，所以 \(H_2=0\)，更高次数也为零。

	在 \(F\otimes E\) 中按同一规则选
	\[
		z'=\frac mg\,y\otimes a-\frac ng\,b\otimes x.
	\]
	由于参与 \(z\) 的每个张量都有一个零次因子，交换符号均为 \(+1\)，直接计算得到 \(\tau(z)=-z'\)。所以在这两组指定的循环群生成元下，次数一的交换同态为乘以 \(-1\)。这不否定 Tor 的对称同构，而是说明把其抽象同构写成具体整数生成元时，必须同时说明生成圈的符号。
\end{solution}

\begin{exercise}\label{g2:ex:torus-map-determinant-cup}
	令 \(u,v\in H^1(T^2;\mathbb Z)\) 分别来自两个正向圆周，并令 \(f:T^2\to T^2\) 为连续映射。若
	\[
		f^*u=au+bv,\qquad f^*v=cu+dv,
		\qquad a,b,c,d\in\mathbb Z,
	\]
	计算 \(f^*\) 在 \(H^2(T^2;\mathbb Z)\) 上的作用。由此判断：若一个环面自映射在 \(H^1\) 上可逆，它在 \(H^2\) 上是否可逆？
\end{exercise}
\begin{solution}
	圆周的二次上同调为零，故由拉回的乘法相容性 \(u^2=v^2=0\)。分次交换律给出 \(vu=-uv\)，而\cref{g2:prop:cross-cup-evaluation}说明 \(uv\) 在两个正向圆周的外积上取值 \(1\)，所以它是 \(H^2(T^2;\mathbb Z)\cong\mathbb Z\) 的生成元。于是
	\[
		f^*(uv)=(au+bv)(cu+dv)
		=ad\,uv+bc\,vu=(ad-bc)\,uv.
	\]
	若 \(f^*\) 在 \(H^1\cong\mathbb Z^2\) 上可逆，其整数矩阵的行列式必为 \(\pm1\)，故在 \(H^2\) 上也是同构。特别地，交换两个圆周方向使矩阵行列式为 \(-1\)，因此把二次生成元变号；这一符号与乘积定向的交换一致。
\end{solution}

\begin{exercise}\label{g2:ex:cup-obstruction-two-open-cover}
	设 \(X=U\cup V\)，其中 \(U,V\) 为开子空间，\(R\) 为含幺交换环。证明：若 \(\alpha\in H^p(X;R)\) 在 \(U\) 上限制为零，\(\beta\in H^q(X;R)\) 在 \(V\) 上限制为零，则 \(\alpha\smile\beta=0\)。据此证明，对 \(p,q\ge1\)，\(S^p\times S^q\) 不存在这样的两开集覆盖：两个包含映射 \(U\to S^p\times S^q\)、\(V\to S^p\times S^q\) 都同伦于常值映射。
\end{exercise}
\begin{solution}
	先设 \(p>0\)，取代表上圈 \(\varphi\)。限制类为零意味着 \(\varphi|_{C_*(U)}=\delta h\)，其中 \(h\) 为 \(U\) 上的 \(p-1\) 次上链。把 \(h\) 在 \(U\) 之外的奇异单形基元上取值为零，就得到 \(X\) 上的延拓 \(\widehat h\)。于是
	\[
		\varphi'=\varphi-\delta\widehat h
	\]
	仍代表 \(\alpha\)，并且在 \(U\) 中的全部链上为零；因此 \(\alpha\) 来自 \(H^p(X,U;R)\)。若 \(p=0\)，零次没有上边界，限制类为零已经说明代表上圈本身在 \(U\) 上为零，仍得同一结论。同理，\(\beta\) 来自 \(H^q(X,V;R)\)。由\cref{g2:prop:relative-cup-product}，两相对提升的杯积属于
	\[
		H^{p+q}(X,U\cup V;R)=H^{p+q}(X,X;R)=0.
	\]
	映到绝对上同调即得 \(\alpha\smile\beta=0\)。

	现在取整数系数，令 \(\alpha,\beta\) 分别为两个球面正向生成元的拉回。它们的杯积在球面基本类的外积上取值 \(1\)，所以非零。若 \(U,V\) 的包含映射都同伦于常值，则任意正次上同调类在它们上的限制都为零：常值映射经一点分解，而一点的正次上同调为零。上一部分将迫使 \(\alpha\smile\beta=0\)，矛盾。
\end{solution}

\begin{exercise}\label{g2:ex:complex-local-orientation}
	在 \(\mathbb C^n\) 上采用实坐标次序 \((x_1,y_1,\ldots,x_n,y_n)\)，令 \(u\) 为 \(H^{2n}(\mathbb C^n,\mathbb C^n\setminus\{0\};\mathbb Z)\) 的正向生成元。计算下列映射对 \(u\) 的拉回：
	\[
		c_r(z_1,\ldots,z_n)
		=(\overline z_1,\ldots,\overline z_r,z_{r+1},\ldots,z_n),
		\qquad 0\le r\le n,
	\]
	以及任意复坐标置换。要求用局部生成元的外积说明符号。
\end{exercise}
\begin{solution}
	先在实轴的局部对 \((\mathbb R,\mathbb R\setminus\{0\})\) 中考察反射。正向小区间的连接边界是“正侧点减负侧点”；反射交换两侧，所以把局部一维同调生成元变号，也把其取值对偶变号。复共轭在一个坐标复平面中保持 \(x\) 不变、把 \(y\) 变号。\cref{g2:prop:local-cross-normalization}把局部二次生成元写成 \(x,y\) 的一维局部生成元的外积，故复共轭把它乘以 \(-1\)。

	再次使用该命题，\(\mathbb C^n\) 的 \(u\) 是 \(n\) 个局部二次生成元按给定坐标次序的外积。外积自然性于是给出
	\[
		c_r^*u=(-1)^r u.
	\]
	交换两个复坐标只交换两个次数为 \(2\) 的局部类，系数为 \((-1)^{2\cdot2}=1\)。任何复坐标置换都可写成相邻交换的复合，所以都保持 \(u\)。这说明一个复坐标块的共轭与两个完整复坐标块的交换，对实定向有不同作用。
\end{solution}

\begin{exercise}\label{g2:ex:surface-ring-genus-bound}
	设 \(g,h\geq1\)，\(f:\Sigma_g\to\Sigma_h\) 连续。选定两曲面的定向，并用二次生成元定义整数 \(d\)：\(f^*\omega_h=d\omega_g\)。证明：若 \(d\ne0\)，则 \(g\geq h\)。此外，若 \(g=h\)，设一次拉回在正文的有序基 \((a_1,b_1,\ldots,a_g,b_g)\) 下的矩阵为 \(P\)，证明 \((\det P)^2=d^{2g}\)。
\end{exercise}
\begin{solution}
	在一次有理上同调上定义双线性配对 \(\lambda_g\)：若 \(uv=r\omega_g\)，令 \(\lambda_g(u,v)=r\)。由\eqref{g2:eq:orientable-surface-products}，它在所列基下的矩阵是 \(g\) 个块 \(\left(\begin{smallmatrix}0&1\\-1&0\end{smallmatrix}\right)\) 的对角和，记为 \(J_g\)。该矩阵可逆。杯积自然性给出
	\[
		\lambda_g(f^*u,f^*v)=d\lambda_h(u,v),\qquad P^{\mathsf T}J_gP=dJ_h.
	\]
	这里 \(P\) 一般为 \(2g\times2h\) 矩阵。若 \(f^*u=0\)，上述等式说明 \(d\lambda_h(u,v)=0\) 对所有 \(v\) 成立。因为 \(d\ne0\) 且 \(J_h\) 可逆，必有 \(u=0\)。所以 \(f^*\) 在一次有理上同调上单射，\(2h\leq2g\)。若 \(g=h\)，则 \(P\) 是方阵，对矩阵恒等式取行列式，利用 \(\det J_g=1\)，得到 \((\det P)^2=d^{2g}\)。
\end{solution}

\begin{exercise}\label{g2:ex:projective-product-even-degree}
	令 \(X=\mathbb CP^1\times\mathbb CP^1\)，两因子的二次生成元拉回后分别记为 \(x,y\)，\(\mathbb CP^2\) 的二次生成元记为 \(u\)。设 \(f:X\to\mathbb CP^2\) 连续且 \(f^*u=ax+by\)，其中 \(a,b\in\mathbb Z\)。计算 \(f^*(u^2)\)，并证明不存在使 \(f^*:H^4(\mathbb CP^2;\mathbb Z)\to H^4(X;\mathbb Z)\) 为同构的连续映射。然后证明任意连续映射 \(r:\mathbb CP^2\times\mathbb CP^2\to\mathbb CP^1\) 在二次整数上同调上的拉回都是零。
\end{exercise}
\begin{solution}
	\cref{g2:thm:complex-projective-ring}给出 \(x^2=y^2=0\)，而外积说明 \(xy\) 是 \(H^4(X;\mathbb Z)\) 的生成元。因 \(x,y\) 的次数均为 \(2\)，交换时没有负号，故
	\[
		f^*(u^2)=(ax+by)^2=2ab\,xy.
	\]
	这是乘以偶数的同态，不能为乘以 \(1\) 或 \(-1\)，所以不可能是整数群的同构。

	对于第二问，把源的两个二次生成元记为 \(v,w\)，靶的二次生成元记为 \(s\)。若 \(r^*s=cv+ew\)，由 \(s^2=0\) 得
	\[
		0=(cv+ew)^2=c^2v^2+2ce\,vw+e^2w^2.
	\]
	整数 Künneth 同构及\cref{g2:thm:complex-projective-ring}说明 \(v^2,vw,w^2\) 是四次上同调的一组自由基。因此 \(c^2=e^2=0\)，从而 \(c=e=0\)。这个障碍来自源与靶的平方关系不同；加性群中的任意同态未必能来自连续映射。
\end{solution}

\begin{exercise}\label{g2:ex:nonorientable-square-kernel}
	在 \(N_h\) 的一次模 \(2\) 上同调中令
	\[
		V=\{v\in H^1(N_h;\mathbb F_2):v^2=0\}.
	\]
	证明 \(V\) 是维数 \(h-1\) 的子空间。若用 \(vw=\lambda(v,w)\omega\) 定义双线性配对，计算它限制在 \(V\) 上的根空间 \(\{v\in V:\lambda(v,w)=0\text{ 对所有 }w\in V\}\)，并区分 \(h\) 的奇偶性。
\end{exercise}
\begin{solution}
	写 \(v=\sum_i a_ic_i\)，\(a_i\in\mathbb F_2\)。由\eqref{g2:eq:nonorientable-surface-products}，有
	\[
		v^2=\Bigl(\sum_i a_i\Bigr)\omega,\qquad
		\lambda\Bigl(\sum_i a_ic_i,\sum_i b_ic_i\Bigr)=\sum_i a_ib_i.
	\]
	故 \(V\) 是系数和为零的超平面，维数为 \(h-1\)。令 \(e=c_1+\cdots+c_h\)，则 \(V=e^\perp\)，其中环境配对就是标准点积，因而非退化。直接验证也可确定 \(V^\perp\)：当 \(h\geq2\) 时，向量 \(c_i+c_j\) 都属于 \(V\)，与它们正交要求各坐标相等，所以 \(V^\perp=\mathbb F_2e\)；当 \(h=1\) 时此结论同样成立。因此限制配对的根空间为 \(V\cap\mathbb F_2e\)。若 \(h\) 为奇数，\(e\) 的系数和为 \(1\)，根空间为零；若 \(h\) 为偶数，\(e\in V\)，根空间就是一维空间 \(\mathbb F_2e\)。
\end{solution}

\begin{exercise}\label{g2:ex:projective-product-ring-automorphisms}
	设 \(1\leq n\leq m\)，令 \(X=\mathbb CP^n\times\mathbb CP^m\)，把两个因子的二次生成元在 \(X\) 上的拉回记为 \(u,v\)。求 \(H^*(X;\mathbb Z)\) 的全部保次数、保单位元的环自同构，区分 \(n<m\) 与 \(n=m\)。证明每个这样的环自同构都由 \(X\) 的某个同胚诱导。
\end{exercise}
\begin{solution}
	由\cref{g2:thm:complex-projective-ring}及无挠的上同调外积，
	\[
		H^*(X;\mathbb Z)=\mathbb Z[u,v]/(u^{n+1},v^{m+1}),\qquad |u|=|v|=2.
	\]
	其加性基为 \(u^iv^j\)、\(0\leq i\leq n\)、\(0\leq j\leq m\)。一个分次环自同构 \(\Phi\) 必须满足
	\[
		\Phi(u)=au+bv,\qquad\Phi(v)=cu+dv,\qquad ad-bc=\pm1.
	\]
	若 \(n<m\)，考察关系 \(\Phi(u)^{n+1}=0\) 中的 \(v^{n+1}\) 系数，它等于 \(b^{n+1}\)。该单项式属于加性基，故 \(b=0\)。行列式条件随即给出 \(a,d\in\{1,-1\}\)。再看 \(\Phi(v)^{m+1}=0\) 中的 \(uv^m\) 系数，得到 \((m+1)cd^m=0\)，所以 \(c=0\)。因此自同构恰是两个生成元各自独立变号。

	若 \(n=m\)，关系 \(\Phi(u)^{n+1}=0\) 的 \(u^nv\) 系数给出 \((n+1)a^nb=0\)，所以 \(a=0\) 或 \(b=0\)。同理，\(c=0\) 或 \(d=0\)。结合行列式为 \(\pm1\)，矩阵只能是两个对角元为 \(\pm1\) 的对角矩阵，或两个反对角元为 \(\pm1\) 的反对角矩阵。也就是说，除了独立变号，还可以交换 \(u,v\)。这些候选确实都保持截断关系，因而都是环自同构。

	最后说明几何实现。复共轭 \([z_0:\cdots:z_r]\mapsto[\overline z_0:\cdots:\overline z_r]\) 是 \(\mathbb CP^r\) 的同胚。在标准 \(\mathbb CP^1\cong S^2\) 上，它在仿射图 \(z=x+\mathrm{i}y\) 中是 \((x,y)\mapsto(x,-y)\)，经球面的立体投影就是一个反射，故由\cref{g2:prop:reflection-antipodal-degree}诱导二次生成元变号。标准包含在二次上同调上为同构，所以复共轭在任意 \(\mathbb CP^r\) 上也使二次生成元变号。对两个因子分别选恒等或复共轭，即实现全部独立变号；当 \(n=m\) 时，再用交换两个因子的同胚实现交换。因而上面列出的每个环自同构都有所要求的实现。
\end{solution}
